% ==========================================================================
% FungalMorphoSpace — SoftwareX manuscript draft
% Template: Elsevier elsarticle (SoftwareX)
% ==========================================================================
\documentclass[review,3p]{elsarticle}

% --- Packages ---
\usepackage{amsmath,amssymb}
\usepackage{graphicx}
\usepackage{hyperref}
\usepackage{xcolor}
\usepackage{booktabs}
\usepackage{siunitx}
\usepackage{listings}
\usepackage{algorithm}
\usepackage{algpseudocode}
\usepackage{lineno}

\modulolinenumbers[5]

% Code listing style
\lstset{
  language=Python,
  basicstyle=\small\ttfamily,
  keywordstyle=\color{blue},
  commentstyle=\color{gray},
  stringstyle=\color{red!70!black},
  numbers=left,
  numberstyle=\tiny\color{gray},
  frame=single,
  breaklines=true,
  captionpos=b,
}

\journal{SoftwareX}

% ==========================================================================
\begin{document}
% ==========================================================================

\begin{frontmatter}

\title{FungalMorphoSpace: A reaction-diffusion simulator for exploring\\
fungal hymenophore morphogenesis}

\author[inst1]{Mario Ahumada Dur\'{a}n}
\ead{mario.ahumada@example.com} % TODO: replace with real email

\affiliation[inst1]{organization={Universidad de Chile},  % TODO: verify
             city={Santiago},
             country={Chile}}

% ------------------------------------------------------------------
\begin{abstract}
FungalMorphoSpace is an open-source Python package for simulating and
quantifying fungal hymenophore patterns using Gierer--Meinhardt
reaction-diffusion systems. The software implements a forward-Euler
solver with CFL-adaptive time-stepping, periodic boundary conditions,
and robust morphometric analysis including dual-metric wavelength
estimation (FFT radial spectrum and autocorrelation, both with quality
control), Euler characteristic computation, and connected-component
spacing statistics. It ships with empirically calibrated parameter sets
for three polypore species---\textit{Fomes fomentarius},
\textit{Lentinus brumalis}, and \textit{Polyporus squamosus}---spanning
an 8$\times$ wavelength range (\SIrange{33}{235}{px},
\SIrange{261}{2043}{\micro\metre}). Key features include temporal
rescaling control ($T_\text{target}$) for cross-species comparability,
a contract-enforced output schema for reproducibility, sensitivity
analysis tools, and an optional three-node (immobile-node) module with a
linear-stability utility that lets users reach the same morphospace at
equal diffusion---making explicit that pattern scale is governed by a
reaction--diffusion length rather than by an extreme diffusion ratio.
Parameters are calibrated empirically, so FungalMorphoSpace is a
hypothesis-generation and morphospace-exploration tool---not an
\emph{ab initio} predictor---for researchers in fungal developmental
biology, computational morphogenesis, and bio-inspired materials design.
\end{abstract}

\begin{keyword}
Reaction-diffusion \sep
Turing patterns \sep
Fungal morphogenesis \sep
Hymenophore \sep
Pattern formation \sep
Gierer--Meinhardt \sep
Polypores \sep
Computational biology
\end{keyword}

\end{frontmatter}

% ------------------------------------------------------------------
\linenumbers

% ==================================================================
\section{Motivation and significance}
\label{sec:motivation}
% ==================================================================

The spore-bearing surfaces of polypore fungi---their
hymenophores---display a remarkable diversity of spatial patterns:
regular hexagonal pores in \textit{Fomes fomentarius}, fine-grained
pores in \textit{Lentinus brumalis}, and coarse pores in
\textit{Polyporus squamosus}. Despite apparent morphological
differences, Kuhar et al.~\cite{kuhar2022} proposed that these
structures represent topological variants of a single self-organizing
mechanism governed by reaction-diffusion (RD) dynamics, as originally
formulated by Turing~\cite{turing1952}.

Testing this hypothesis computationally requires tools that can
(i)~generate patterns across diverse parameter regimes, (ii)~quantify
morphological features objectively, and (iii)~compare simulated
outputs with empirical measurements. While general-purpose RD
simulators exist, none provides an integrated pipeline tailored to
fungal morphogenesis that includes species-specific calibration,
topological analysis, and reproducible output contracts.

FungalMorphoSpace addresses this gap by providing:
\begin{itemize}
  \item A modular Gierer--Meinhardt simulator with CFL-adaptive
        stability and selectable kinetics;
  \item Quantitative morphometric and topological analysis
        (wavelength, Euler characteristic, spacing statistics);
  \item Calibrated parameter presets for three validated polypore
        species;
  \item Reproducible, contract-enforced output schemas;
  \item Sensitivity analysis tools for systematic morphospace
        exploration.
\end{itemize}

The software targets three user communities: (a)~mycologists and
developmental biologists studying pattern formation in fungi,
(b)~computational biologists modeling self-organization, and
(c)~materials scientists exploring bio-inspired porous architectures.


% ==================================================================
\section{Software description}
\label{sec:description}
% ==================================================================

% ------------------------------------------------------------------
\subsection{Software architecture}
\label{sec:architecture}

FungalMorphoSpace is organized into five modules
(Fig.~\ref{fig:architecture}):

\begin{description}
  \item[\texttt{core}] Turing pattern simulator with polymorphic
    kinetics models (Gierer--Meinhardt, Schnakenberg, Gray--Scott).
  \item[\texttt{analysis}] Topology analyzer (Euler characteristic,
    connected components) and morphometric tools (FFT and
    autocorrelation wavelength, spacing statistics, isotropy index).
  \item[\texttt{runners}] Validation pipeline orchestrating
    single-species, multi-species, and sensitivity runs with progress
    reporting and statistical aggregation.
  \item[\texttt{contracts}] Output schema definitions ensuring
    machine-readable CSV reproducibility across versions.
  \item[\texttt{utils}] Parameter sweep tools and repository cleanup
    utilities.
\end{description}

% TODO: Generate architecture diagram and place in figures/
% \begin{figure}[ht]
%   \centering
%   \includegraphics[width=0.8\textwidth]{figures/architecture.pdf}
%   \caption{Modular architecture of FungalMorphoSpace.}
%   \label{fig:architecture}
% \end{figure}

% ------------------------------------------------------------------
\subsection{Mathematical model}
\label{sec:model}

The core simulator implements the Gierer--Meinhardt
system~\cite{gierer1972}:

\begin{equation}
\label{eq:gm}
\begin{aligned}
  \frac{\partial u}{\partial t}
    &= D_u \nabla^2 u + \rho\!\left(\frac{u^2}{v + \varepsilon}
      - b\,u + a\right), \\[4pt]
  \frac{\partial v}{\partial t}
    &= D_v \nabla^2 v + \rho\!\left(u^2 - v\right),
\end{aligned}
\end{equation}

\noindent where $u$ is a short-range activator, $v$ a long-range
inhibitor, $D_u$ and $D_v$ are diffusion coefficients, $\rho$ is
the kinetic intensity (a proxy for metabolic activity), and
$\varepsilon = 10^{-6}$ prevents numerical singularity in the
autocatalytic term.

The homogeneous steady state is
$u_0 = (a + 1)/b$, $v_0 = u_0^2$. Classical two-species Turing
instability requires a diffusion disparity $D_v/D_u \gg 1$, and the
pattern wavelength then depends jointly on $D_v/D_u$ and $\rho$. This
disparity is not fundamental, however: it is one way to set the
underlying reaction--diffusion length $\ell \sim \sqrt{D_\text{eff}/k}$
that fixes the wavelength. An immobile (non-diffusing) third species
coupled to the inhibitor produces the same instability at \emph{equal}
diffusion, as established for binding substrates and morphogen
adsorption~\cite{korvasova2015,nesterenko2017} and catalogued for
three-node networks with equally diffusing signals~\cite{marcon2016}.
FungalMorphoSpace includes an optional three-node module
(Sec.~\ref{sec:threenode}) implementing this mechanism, which lets users
contrast disparity-driven and immobile-node routes to the same morphospace.

\subsection{Three-node (immobile-node) module}
\label{sec:threenode}

The module \texttt{core.three\_node} adds two diffusing morphogens and a
third immobile field $w$ (interpreted as local hyphal density), with a
linear-stability utility (\texttt{dispersion\_growth\_rate},
\texttt{fastest\_growing\_wavelength}) that computes the Turing band
analytically from the reaction Jacobian. It is intended for exploring
\emph{whether} a given topology patterns and at what scale \emph{before}
committing to a simulation. We stress that the underlying equal-diffusion
mechanism is established prior art~\cite{korvasova2015,nesterenko2017,marcon2016,raspopovic2014};
the module's role here is to make that mechanism convenient to apply to
hymenophore morphospace, not to introduce a new mechanism.

\paragraph{Temporal rescaling.}
Under the change of variable $\tau = \rho\,t$, the effective
diffusion coefficients become $D_u^* = D_u / \rho$ and
$D_v^* = D_v / \rho$. This reveals that the characteristic wavelength
$\lambda$ depends on the ratio $D_v / \rho$, not $D_v$ alone---a key
insight that resolves cases where increasing $D_v/D_u$ paradoxically
\emph{decreases} $\lambda$ when $\rho$ increases simultaneously
(the ``Brumalis paradox''; see Table~\ref{tab:species}).

% ------------------------------------------------------------------
\subsection{Numerical implementation}
\label{sec:numerical}

Spatial discretization uses a second-order five-point finite-difference
stencil on a square grid ($N \times N$, $N \in \{256, 512, 1024\}$)
with periodic boundary conditions. Time integration uses the explicit
(forward) Euler method with adaptive time-stepping via the CFL
condition:

\begin{equation}
\label{eq:cfl}
  \Delta t \leq \frac{(\Delta x)^2}{4\,\max(D_u,\, D_v)},
\end{equation}

\noindent reduced by a safety factor of $0.8$ to ensure stability for
stiff reactions at high $D_v/D_u$ ratios. Temporal comparability
across species is enforced through a target physical time $T_\text{target}$,
from which the number of steps is computed as
$n_\text{steps} = \lceil T_\text{target} / \Delta t_\text{final}\rceil$.

A critical numerical fix prevents catastrophic blowup in the
$u^2 / v$ term: a floor value $v_\text{floor} = 10^{-6}$ is enforced
after each time step, complementing the $\varepsilon$ regularization
in the kinetic function.

% ------------------------------------------------------------------
\subsection{Morphometric and topological analysis}
\label{sec:analysis}

Pattern quantification employs two complementary approaches:

\paragraph{Wavelength estimation.}
Two independent methods provide robustness:
\begin{enumerate}
  \item \textbf{FFT radial spectrum}: the 2D power spectrum is
    radially averaged; the dominant peak above a minimum frequency
    ($\geq 2$ cycles per domain) yields $\lambda_\text{FFT}$.
    Quality control (QC) requires peak prominence
    $\geq 3\times$ median background.
  \item \textbf{Autocorrelation}: the 2D autocorrelation (via the
    Wiener--Khinchin theorem) is radially averaged; $\lambda_\text{AC}$
    is the position of the first maximum after the first minimum.
\end{enumerate}

The pipeline selects the best estimate via a QC cascade:
FFT (if QC pass) $\to$ autocorrelation (if QC pass) $\to$ fallback
(first finite positive value, marked as QC fail).

\paragraph{Topological metrics.}
\begin{itemize}
  \item \textbf{Euler characteristic}
    $\chi = N_\text{components} - N_\text{holes}$, computed via
    \texttt{skimage.measure.euler\_number}~\cite{euler_number_skimage}
    with 8-connectivity. Spots yield $\chi > 0$; labyrinths
    $\chi \approx 0$; networks $\chi < 0$.
  \item \textbf{Connected components}: count and centroids of
    binarized foreground regions, replacing the noise-sensitive local
    maxima detection used in earlier versions.
  \item \textbf{Spacing statistics}: nearest-neighbor distances from
    component centroids (via KD-tree), yielding mean spacing and
    coefficient of variation as regularity measures.
  \item \textbf{Isotropy index}: angular distribution of 2D FFT power;
    $I = 1$ for perfectly isotropic patterns.
\end{itemize}

% ------------------------------------------------------------------
\subsection{Calibrated species database}
\label{sec:species}

Table~\ref{tab:species} summarizes the three validated species.
Parameters were calibrated empirically by adjusting $D_v/D_u$,
$\rho$, and $b$ until simulated $\lambda$ matched observed pore
spacings from the literature. The spatial calibration factor
$\SI{8.7}{\micro\metre/px}$ derives from \textit{F.~fomentarius}
scanning electron microscopy data~\cite{klemm2024}.

\begin{table}[ht]
\centering
\caption{Calibrated species parameters and validation results.}
\label{tab:species}
\begin{tabular}{@{}lccccccc@{}}
\toprule
Species & Key & $D_v/D_u$ & $\rho$ & $b$ & Grid
  & $\lambda$ (px) & $\lambda$ (\si{\micro\metre}) \\
\midrule
\textit{F.~fomentarius} & \texttt{fomes}
  & 150 & 0.2 & 1.0 & 512 & 46 & 400 \\
\textit{L.~brumalis}    & \texttt{brumalis}
  & 250 & 0.4 & 1.0 & 512 & 33 & 261 \\
\textit{P.~squamosus}   & \texttt{squamosus}
  & 3750 & 5.0 & 3.0 & 1024 & 235 & 2043 \\
\bottomrule
\end{tabular}
\end{table}

\paragraph{Transparency note.}
These parameters were obtained by empirical fitting, not from
independent biophysical measurements. The software therefore serves as
an exploration and hypothesis-generation tool rather than an
\emph{ab initio} predictor. Strategies for non-tautological parameter
anchoring---e.g., deriving $D_\text{eff}$ from measured substrate
porosity---are discussed in the project documentation.


% ==================================================================
\section{Illustrative examples}
\label{sec:examples}
% ==================================================================

% ------------------------------------------------------------------
\subsection{Multi-species validation}

\begin{lstlisting}[caption={Full validation suite.},label=lst:validation]
# Validate all 3 species (5 replicates, high resolution)
python scripts/run_integrated_validation.py \
    --species all --n_runs 5 --grid 1024
\end{lstlisting}

This generates a comprehensive comparison figure with pattern images,
wavelength measurements, and a scaling-law plot. Across the three
species, log-log regression yields
$\lambda_\text{sim} \propto \text{Size}_\text{bio}^{0.96}$
($R^2 \approx 0.997$, $N = 3$; illustrative, not statistically
powered).

% ------------------------------------------------------------------
\subsection{Programmatic API}

\begin{lstlisting}[caption={Python API usage.},label=lst:api]
from fungalmorphospace.core import (
    TuringSimulator, GiererMeinhardtKinetics)
from fungalmorphospace.analysis import TopologyAnalyzer

# Configure and run
kinetics = GiererMeinhardtKinetics(
    rho=0.2, a=0.1, b=1.0)
sim = TuringSimulator(
    kinetics_model=kinetics,
    D_u=0.1, D_v=15.0, grid_size=512)
sim.initialize(perturbation_amplitude=0.1)
sim.run(num_steps=10000)

# Analyze
analyzer = TopologyAnalyzer(sim.u, dx=1.0)
metrics = analyzer.compute_all_metrics()
print(f"Lambda: {metrics['wavelength_fft']:.1f} px")
print(f"Euler: {metrics['euler_characteristic']}")
\end{lstlisting}

% ------------------------------------------------------------------
\subsection{Topological continuity}

The script \texttt{test\_laminillas.py} demonstrates that pores and
gills are topological variants: varying $D_v/D_u$ while tracking
$\chi$ reveals a smooth transition from isolated spots ($\chi > 0$)
through labyrinths ($\chi \approx 0$) to connected networks
($\chi < 0$), all generated by the same Gierer--Meinhardt mechanism.

% ------------------------------------------------------------------
\subsection{Sensitivity analysis}

\begin{lstlisting}[caption={Parameter sweep.},label=lst:sweep]
from fungalmorphospace.utils.sensitivity_analysis \
    import SensitivityAnalyzer

sa = SensitivityAnalyzer(base_params={
    'D_u': 1.0, 'grid_size': 256,
    'dx': 1.0, 'dt': 0.0005,
    'steps': 5000, 'model': 'gierer_meinhardt'})

df = sa.sweep_2d(
    'D_v_D_u', [50, 100, 200, 400],
    'rho', [0.1, 0.2, 0.5, 1.0])
\end{lstlisting}

% ------------------------------------------------------------------
\subsection{Equal-diffusion exploration and the wavelength scale law}
\label{sec:equaldiff}

The three-node module (Sec.~\ref{sec:threenode}) lets users probe whether
the hymenophore morphospace can be reached \emph{without} a diffusion
disparity. Using the analytical dispersion utility, a topology in which
the immobile node couples to the inhibitor is Turing-unstable at
$D_u = D_v$, and its fastest-growing wavelength follows the
reaction--diffusion length scaling $\lambda \propto \sqrt{D}$
(Fig.~\ref{fig:scalelaw}). Nonlinear simulations reproduce this scaling:
the measured spot spacing tracks $\sqrt{D}$ with a constant offset above
the linear $\lambda^\ast$ (coarsening factor $\approx 2.4$, independent of
$D$ and of the saturation strength). The scripts
\texttt{search\_turing\_topology.py}, \texttt{confirm\_three\_node.py},
and \texttt{figure\_scale\_law.py} reproduce these results.

This example illustrates the software's analytical value: it shows that the
extreme ratio $D_v/D_u = 3750$ used for the coarse-pored
\textit{P.~squamosus} preset (Table~\ref{tab:species}) is not a necessary
feature of the model but one parameterization of the diffusion length;
the same scale is reachable at equal diffusion via the immobile node.

\begin{figure}[ht]
  \centering
  \includegraphics[width=0.78\textwidth]{../results/three_node/figure_scale_law.png}
  \caption{Wavelength versus $\sqrt{D}$ at equal diffusion ($D_u=D_v$) for
  the immobile-node topology. Linear dispersion prediction (circles) and
  nonlinear simulated spacing (squares) are both straight lines in
  $\sqrt{D}$; simulated points (stars) lie on the nonlinear curve. Dotted
  lines mark the biological pore spacings of the three calibrated species.
  Generated by \texttt{figure\_scale\_law.py}.}
  \label{fig:scalelaw}
\end{figure}


% ==================================================================
\section{Impact}
\label{sec:impact}
% ==================================================================

\subsection{Scientific contributions}

FungalMorphoSpace has supported:
\begin{itemize}
  \item Computational \emph{exploration} of the Kuhar
    hypothesis~\cite{kuhar2022}: the three calibrated species, spanning an
    8$\times$ wavelength range, are reproducible as parametric variants of
    one Gierer--Meinhardt mechanism (consistency, not independent
    validation; see Limitations);
  \item A transparent treatment of the $D/\rho$ nondimensionalization
    (effective diffusion), which clarifies cases where increasing
    $D_v/D_u$ while raising $\rho$ \emph{decreases} $\lambda$;
  \item Demonstration of topological continuity between morphotypes
    (spots $\leftrightarrow$ labyrinths) within a single model;
  \item An analytical demonstration that the morphospace scale is governed
    by a diffusion length and is reachable at equal diffusion via an
    immobile node (Sec.~\ref{sec:equaldiff}), so that extreme
    diffusion ratios are not a necessary model ingredient.
\end{itemize}

\subsection{Broader applications}

Beyond mycology, the software provides:
\begin{itemize}
  \item A pedagogical tool for teaching pattern formation and
    self-organization in undergraduate or graduate courses;
  \item A generative engine for bio-inspired porous architectures
    in biomaterials design;
  \item A baseline framework for integrating molecular data with
    morphogenetic models.
\end{itemize}

\subsection{Limitations and future directions}

\textbf{Current limitations:}
\begin{itemize}
  \item \emph{Empirical calibration (the central caveat).} Parameters are
    fitted to observed $\lambda$, not derived from independent biophysical
    measurements, so reproducing a wavelength is a consistency check, not a
    prediction. Because $\lambda$ is set by a single effective
    reaction--diffusion length, many parameter combinations and indeed
    several distinct mechanisms reproduce the same $\lambda$; breaking this
    degeneracy requires measuring the underlying diffusion and reaction
    rates (Future directions).
  \item \emph{The pattern-forming mechanism is not novel.} Equal-diffusion
    Turing via an immobile/binding species is established theory, including
    applications to other organisms~\cite{korvasova2015,nesterenko2017,marcon2016,raspopovic2014};
    the contribution here is the application to fungal hymenophores, the
    accompanying reproducible tooling, and the demonstration that the prior
    extreme-ratio parameterization is unnecessary.
  \item \emph{Stability caveat for the immobile-node module.}
    Reaction-diffusion-ODE systems with an autocatalytic immobile component
    need not admit stable Turing patterns~\cite{marciniakczochra2016}; the
    preset topology used here is non-autocatalytic (binding-type), but
    long-time stability of the nonlinear patterns should be verified per
    application.
  \item Simulations are 2D; real hymenophores are 3D surfaces
    (locally planar approximations are reasonable).
  \item Only $N = 3$ species are calibrated, giving limited statistical
    power for scaling-law inference.
  \item The explicit Euler solver, while stable with CFL control, becomes
    expensive at very high $D_v/D_u$ ratios and large grids
    (e.g.\ \textit{P.~squamosus}-scale wavelengths).
\end{itemize}

\textbf{Planned extensions:}
\begin{itemize}
  \item Non-tautological parameter anchoring via independently measured
    diffusion and reaction rates---e.g.\ substrate porosity and
    microstructure~\cite{klemm2024} and \mbox{Ca\textsuperscript{2+}}
    signal-propagation data in mycelial networks~\cite{itani2023}---which
    would discriminate the candidate mechanisms rather than fit the
    wavelength;
  \item Extension to $N > 10$ species for robust allometric scaling;
  \item 3D simulations with surface curvature effects;
  \item Machine learning for parameter inference from microscopy
    images;
  \item Implicit or semi-implicit solvers for improved efficiency.
\end{itemize}


% ==================================================================
\section{Conclusions}
\label{sec:conclusions}
% ==================================================================

FungalMorphoSpace demonstrates that the morphological diversity of
fungal hymenophores is consistent with parametric variation within a
single Gierer--Meinhardt reaction-diffusion framework. The software
provides a reproducible, well-documented pipeline for simulating
patterns, extracting topological and morphometric features, and
exploring the accessible morphospace. While the current empirical
calibration approach limits \emph{ab initio} prediction, the tool
enables systematic hypothesis generation and complements experimental
studies of fungal developmental biology.

The package is available under a CC~BY-NC~4.0 license at
\url{https://github.com/mahumadad/FungalMorphoSpace}.


% ==================================================================
\section*{Code and data availability}
% ==================================================================
The source code is available under a CC~BY-NC~4.0 license at
\url{https://github.com/mahumadad/FungalMorphoSpace}. A versioned release
will be archived with a DOI (e.g.\ Zenodo) upon publication. All figures and
tables in this paper are reproducible with the scripts under
\texttt{scripts/} and the configuration in \texttt{data/species\_data.json}.

% ==================================================================
\section*{CRediT author contributions}
% ==================================================================
\textbf{Mario Ahumada Dur\'{a}n:} Conceptualization, Methodology, Software,
Validation, Investigation, Writing -- original draft, Writing -- review \&
editing.

% ==================================================================
\section*{Declaration of generative AI use}
% ==================================================================
The author used an AI assistant (Claude) to support code review, literature
search and verification, and drafting. All AI-assisted outputs---including
every cited reference---were verified by the author, who takes full
responsibility for the content.

% ==================================================================
\section*{Funding}
% ==================================================================
% TODO (author): state funding sources, or "This research received no
% specific grant from funding agencies in the public, commercial, or
% not-for-profit sectors."

% ==================================================================
\section*{Conflict of interest}
% ==================================================================
The author declares no competing financial or non-financial interests.

% ==================================================================
\section*{Acknowledgements}
% ==================================================================
% TODO (author): add acknowledgements as appropriate.


% ==================================================================
\bibliographystyle{elsarticle-num}
\bibliography{references}
% ==================================================================

\end{document}
